% GENERATED FILE. Edit canonical lessons, then run npm run generate:books.
% canonical-source-hash: fnv1a64:175ea21df6d10081
% canonical-lessons: IT-C482-sopracciglio, IT-C482-ciglia, IT-C482-palpebra, IT-C482-tempia, IT-C482-mascella

\chapter{Il viso da vicino}
\label{ch:it-il-viso-da-vicino}
\hlchaptermodality{482}

\begin{chapteropening}
\textbf{By the end of this chapter:} \emph{I can describe the face in detail: brows, lashes, eyelids, temples and jaw.}

Complete the last lesson of chapter 482: I can describe the face in detail: brows, lashes, eyelids, temples and jaw.
\end{chapteropening}

\section[il sopracciglio]{il sopracciglio — eyebrow}

\label{lesson:IT-C482-sopracciglio}



\begin{quote}
Before the new one: say the Italian for heel, then the Italian for big toe.
\end{quote}

\subsection*{You'll want to know: il sopracciglio}
\textbf{il sopracciglio} — \textquotedblleft{}eyebrow\textquotedblright{}.

\begin{quote}
Quando è sorpreso alza un sopracciglio.
\end{quote}

\begin{quote}
\emph{When he is surprised he raises an eyebrow.}
\end{quote}

An irregular plural: \emph{le sopracciglia}.

\begin{cousinweb}
\emph{sopracciglio} is inherited from Latin \emph{supercilium} \textquotedblleft{}eyebrow\textquotedblright{} (Lewis \& Short), built on \emph{super} \textquotedblleft{}above\textquotedblright{} and \emph{cilium} \textquotedblleft{}eyelid\textquotedblright{}.
\end{cousinweb}

\begin{grammarlens}[title={What you've built}]
One more word for telling, explaining and describing at B1.
\end{grammarlens}

\subsection*{Your turn}
\begin{itemize}
\raggedright
  \item \emph{Say it:} \emph{il sopracciglio}
  \item \emph{Say it:} \emph{il sopracciglio}, once more
  \item \emph{Say it:} a sentence of your own with \emph{il sopracciglio}
  \item \emph{Recall:} say the Italian for heel, then the Italian for big toe, then say \emph{il sopracciglio} again
\end{itemize}

\begin{tcolorbox}[breakable,colback=teal!4,colframe=teal!35!black,title={Before you move on}]
What does il sopracciglio mean? (\textquotedblleft{}Eyebrow\textquotedblright{}.) Say it once more.
\end{tcolorbox}
\section[le ciglia]{le ciglia — eyelashes}

\label{lesson:IT-C482-ciglia}



\begin{quote}
Before the new one: say the Italian for big toe, then the Italian for eyebrow.
\end{quote}

\subsection*{You'll want to know: le ciglia}
\textbf{le ciglia} — \textquotedblleft{}eyelashes\textquotedblright{}.

\begin{quote}
La bambina ha le ciglia lunghe.
\end{quote}

\begin{quote}
\emph{The little girl has long eyelashes.}
\end{quote}

The singular \emph{il ciglio} also means an edge: \emph{il ciglio della strada} — the roadside.

\begin{cousinweb}
\emph{ciglia} is inherited from Latin \emph{cilia}, the plural of \emph{cilium} \textquotedblleft{}eyelid\textquotedblright{} (Lewis \& Short).
\end{cousinweb}

\begin{grammarlens}[title={What you've built}]
One more word for telling, explaining and describing at B1.
\end{grammarlens}

\subsection*{Your turn}
\begin{itemize}
\raggedright
  \item \emph{Say it:} \emph{le ciglia}
  \item \emph{Say it:} \emph{le ciglia}, once more
  \item \emph{Say it:} a sentence of your own with \emph{le ciglia}
  \item \emph{Recall:} say the Italian for big toe, then the Italian for eyebrow, then say \emph{le ciglia} again
\end{itemize}

\begin{tcolorbox}[breakable,colback=teal!4,colframe=teal!35!black,title={Before you move on}]
What does le ciglia mean? (\textquotedblleft{}Eyelashes\textquotedblright{}.) Say it once more.
\end{tcolorbox}
\section[la palpebra]{la palpebra — eyelid}

\label{lesson:IT-C482-palpebra}



\begin{quote}
Before the new one: say the Italian for eyebrow, then the Italian for eyelashes.
\end{quote}

\subsection*{You'll want to know: la palpebra}
\textbf{la palpebra} — \textquotedblleft{}eyelid\textquotedblright{}.

\begin{quote}
Ho le palpebre pesanti, ho sonno.
\end{quote}

\begin{quote}
\emph{My eyelids are heavy, I am sleepy.}
\end{quote}

The plural is \emph{le palpebre}.

\begin{cousinweb}
\emph{palpebra} comes from Latin \emph{palpebra} \textquotedblleft{}eyelid\textquotedblright{} (EtymDB), related to \emph{palpāre} \textquotedblleft{}to touch softly\textquotedblright{} (Lewis \& Short).
\end{cousinweb}

\begin{grammarlens}[title={What you've built}]
One more word for telling, explaining and describing at B1.
\end{grammarlens}

\subsection*{Your turn}
\begin{itemize}
\raggedright
  \item \emph{Say it:} \emph{la palpebra}
  \item \emph{Say it:} \emph{la palpebra}, once more
  \item \emph{Say it:} a sentence of your own with \emph{la palpebra}
  \item \emph{Recall:} say the Italian for eyebrow, then the Italian for eyelashes, then say \emph{la palpebra} again
\end{itemize}

\begin{tcolorbox}[breakable,colback=teal!4,colframe=teal!35!black,title={Before you move on}]
What does la palpebra mean? (\textquotedblleft{}Eyelid\textquotedblright{}.) Say it once more.
\end{tcolorbox}
\section[la tempia]{la tempia — temple (side of the head)}

\label{lesson:IT-C482-tempia}



\begin{quote}
Before the new one: say the Italian for eyelashes, then the Italian for eyelid.
\end{quote}

\subsection*{You'll want to know: la tempia}
\textbf{la tempia} — \textquotedblleft{}temple (side of the head)\textquotedblright{}.

\begin{quote}
Ho un dolore alla tempia sinistra.
\end{quote}

\begin{quote}
\emph{I have a pain in my left temple.}
\end{quote}

The plural is \emph{le tempie}.

\begin{cousinweb}
\emph{tempia} comes from Latin \emph{tempora} \textquotedblleft{}the temples of the head\textquotedblright{} (EtymDB; Lewis \& Short), through a Vulgar Latin form \emph{tempula}.
\end{cousinweb}

\begin{grammarlens}[title={What you've built}]
One more word for telling, explaining and describing at B1.
\end{grammarlens}

\subsection*{Your turn}
\begin{itemize}
\raggedright
  \item \emph{Say it:} \emph{la tempia}
  \item \emph{Say it:} \emph{la tempia}, once more
  \item \emph{Say it:} a sentence of your own with \emph{la tempia}
  \item \emph{Recall:} say the Italian for eyelashes, then the Italian for eyelid, then say \emph{la tempia} again
\end{itemize}

\begin{tcolorbox}[breakable,colback=teal!4,colframe=teal!35!black,title={Before you move on}]
What does la tempia mean? (\textquotedblleft{}Temple (side of the head)\textquotedblright{}.) Say it once more.
\end{tcolorbox}
\section[la mascella]{la mascella — jaw}

\label{lesson:IT-C482-mascella}



\begin{quote}
Before the new one: say the Italian for eyelid, then the Italian for temple.
\end{quote}

\subsection*{You'll want to know: la mascella}
\textbf{la mascella} — \textquotedblleft{}jaw\textquotedblright{}.

\begin{quote}
Quando è nervoso stringe la mascella.
\end{quote}

\begin{quote}
\emph{When he is nervous he clenches his jaw.}
\end{quote}

The lower jaw is also \emph{la mandibola}.

\begin{cousinweb}
\emph{mascella} is inherited from Latin \emph{maxilla} \textquotedblleft{}jaw\textquotedblright{} (EtymDB), a little form of \emph{māla} \textquotedblleft{}cheekbone\textquotedblright{} (Lewis \& Short).
\end{cousinweb}

\begin{grammarlens}[title={What you've built}]
One more word for telling, explaining and describing at B1.
\end{grammarlens}

\subsection*{Your turn}
\begin{itemize}
\raggedright
  \item \emph{Say it:} \emph{la mascella}
  \item \emph{Say it:} \emph{la mascella}, once more
  \item \emph{Say it:} a sentence of your own with \emph{la mascella}
  \item \emph{Recall:} say the Italian for eyelid, then the Italian for temple, then say \emph{la mascella} again
\end{itemize}

\begin{tcolorbox}[breakable,colback=teal!4,colframe=teal!35!black,title={Before you move on}]
What does la mascella mean? (\textquotedblleft{}Jaw\textquotedblright{}.) Say it once more.
\end{tcolorbox}
